\indent ಈ ಉದಾಹರಣೆಗಳಲ್ಲಿ ನಾಮಪದಗಳ ಮುಂದೆ “ಅಂತೆ” ಎಂಬ ಪ್ರತ್ಯಯ ಸೇರಿ ಹೊಸ ಶಬ್ದಗಳಾಗಿರುವುದನ್ನು ಕಾಣುತ್ತೇವೆ ; ಈ ಪ್ರತಿ ಯೊಂದು ಶಬ್ದದಲ್ಲೂ ನಾಮಪದದಲ್ಲಿರುವ ವಿಭಕ್ತಿ ಪ್ರತ್ಯಯವು ಹಾಗೆಯೇ ಉಳಿದುಕೊಂಡಿರುವುದು ಗಮನಾರ್ಹ ಅಂಶ.\par

\indent ಹೀಗೆ ನಾಮಪದಗಳ ಮುಂದೆ ಅಂತೆ, ತನಕ, ವರೆಗೆ ಇಂತ, ಆಗಿ ವೋಲ್‌, ಮಟ್ಟಿಗೆ, ಓಸ್ಕರ, ಸಲುವಾಗಿ ಇತ್ಯಾದಿ ಪ್ರತ್ಯಯಗಳು ಸೇರಿ ಆಗುವ ಹೊಸ ಶಬ್ದಗಳಿಗೆ ತದ್ಧಿತಾಂತಾವ್ಯಯಗಳೆಂದು ಹೆಸರು.\par

\vspace{3pt plus 2pt}\noindent\hspace*{1.2em}\kbvlabel{ಉದಾ}\par\vspace{1pt}

\vspace{5pt plus 4pt minus 1pt}
\noindent{\renewcommand{\arraystretch}{1.24}\setlength{\tabcolsep}{4pt}\rowcolors{1}{kbvrow}{white}
\begin{tabularx}{\textwidth}{X X X}
ದೊರೆಯತನಕ, & ತಿಂಗಳವರೆಗೆ, & ಹಣಕ್ಕೋಸ್ಕರ, \\
ರಾಣಿಯತನಕ, & ಕೃಷ್ಣನಿಗಿಂತ, & ನನ್ನಮಟ್ಟಿಗೆ, \\
ಬಾಗಿಲತನಕ, & ಶಿಲೆಗಿಂತ, & ಅವಳಸಲುವಾಗಿ, \\
ಜನಗಳವರೆಗೆ, & ಪುಸ್ತಕಗಳಿಗಿಂತ, & ರಾಜನವೋಲ್‌ \\
ದಿಗಂತದವರೆಗೆ & ಅವರಿಗಾಗಿ & ಸಿಂಹದವೋಲ್‌ \\
\end{tabularx}}
\vspace{5pt plus 4pt minus 1pt}

\kbvexercise{ಅಭ್ಯಾಸ ಪತ್ರಿಕೆ ೧೫}

\vspace{2.5pt plus 3pt minus 0.5pt}\noindent\hspace*{1.2em}\kbvnum{(1)}\hspace{0.55em}ತದ್ಧಿತ ಪ್ರತ್ಯಯ ಎಂದರೇನು? ಒಂದು ಉದಾಹರಣೆ ಸಹಿತ ತಿಳಿಸಿ.\par

\vspace{2.5pt plus 3pt minus 0.5pt}\noindent\hspace*{1.2em}\kbvnum{(2)}\hspace{0.55em}ತದ್ಧಿತಾಂತಗಳೆಂದರೇನು? ಎರಡು ಉದಾಹರಣೆಗಳನ್ನು ಕೊಡಿ.\par

\vspace{2.5pt plus 3pt minus 0.5pt}\noindent\hspace*{1.2em}\kbvnum{(3)}\hspace{0.55em}ತದ್ಧಿತಾಂತಗಳ ಮೂರು ಪ್ರಕಾರಗಳನ್ನು ತಿಳಿಸಿ.\par

\vspace{2.5pt plus 3pt minus 0.5pt}\noindent\hspace*{1.2em}\kbvnum{(4)}\hspace{0.55em}ತದ್ಧಿತಾಂತನಾಮ ಎಂದರೇನು? ಉದಾಹರಣೆ ಕೊಡಿ.\par

\vspace{2.5pt plus 3pt minus 0.5pt}\noindent\hspace*{1.2em}\kbvnum{(5)}\hspace{0.55em}ತದ್ಧಿತಾಂತ ಭಾವನಾಮ ಎಂದರೇನು? ಉದಾಹರಣೆ ಸಹಿತ ವಿವರಿಸಿ.\par

\vspace{2.5pt plus 3pt minus 0.5pt}\noindent\hspace*{1.2em}\kbvnum{(6)}\hspace{0.55em}ಕೃದಂತ ಭಾವನಾಮಗಳಿಗೂ ತದ್ಧಿತಾಂತ ಭಾವನಾಮಗಳಿಗೂ ಇರುವ ವ್ಯತ್ಯಾಸವೇನು?\par

\vspace{2.5pt plus 3pt minus 0.5pt}\noindent\hspace*{1.2em}\kbvnum{(7)}\hspace{0.55em}ಬ್ರಾಜ್ಮಟಕೆ, ನಾಚಿಕೆ-ಈ ಪದಗಳಲ್ಲಿ ಕಂಡುಬರುವ ವ್ಯಾಕರಣ ವಿಶೇಷವೇನು?\par
